%%%%%%%%%%%%%%%%%%%%%%%%%%%%%%%%%%%%%%%%%%%%%%%%%%%%%%%%%%%%%%%%%%%%%%%%%%%
%%  Project: CPHOS LaTeX Template
%%  File: example-problem.tex
%%  Copyright © 2026 gameswu & CPHOS (www.cphos.cn)
%%
%%  Permission is hereby granted, free of charge, to any person obtaining a 
%%  copy of this software and associated documentation files (the "Software"), 
%%  to deal in the Software without restriction, including without limitation 
%%  the rights to use, copy, modify, merge, publish, distribute, sublicense, 
%%  and/or sell copies of the Software, and to permit persons to whom the 
%%  Software is furnished to do so, subject to the following conditions:
%%
%%  The above copyright notice and this permission notice shall be included 
%%  in all copies or substantial portions of the Software.
%%
%%  THE SOFTWARE IS PROVIDED "AS IS", WITHOUT WARRANTY OF ANY KIND, EXPRESS 
%%  OR IMPLIED, INCLUDING BUT NOT LIMITED TO THE WARRANTIES OF MERCHANTABILITY, 
%%  FITNESS FOR A PARTICULAR PURPOSE AND NONINFRINGEMENT. IN NO EVENT SHALL 
%%  THE AUTHORS OR COPYRIGHT HOLDERS BE LIABLE FOR ANY CLAIM, DAMAGES OR OTHER 
%%  LIABILITY, WHETHER IN AN ACTION OF CONTRACT, TORT OR OTHERWISE, ARISING 
%%  FROM, OUT OF OR IN CONNECTION WITH THE SOFTWARE OR THE USE OR OTHER 
%%  DEALINGS IN THE SOFTWARE.
%%%%%%%%%%%%%%%%%%%%%%%%%%%%%%%%%%%%%%%%%%%%%%%%%%%%%%%%%%%%%%%%%%%%%%%%%%%

% CPHOS理论命题模板

\documentclass[answer]{cphos}

% === 设置文档信息 ===
\cphostitle{CPHOS物理竞赛联考}
\cphossubtitle{理论试题命题模板}

\setscorecheck{true}  % 启用分值检查，可设为false关闭

\begin{document}
\begin{problem}[40]{潘宁阱} 

% === 题干 ===
\begin{problemstatement}
潘宁阱是一种利用均匀磁场和四极电场来囚禁带电粒子的装置。它就像一个精准的“粒子牢笼”，可以长时间稳定地束缚单个离子或电子，是实现高精度测量的关键工具。

在潘宁阱内建立直角坐标系，空间中的电势分布为：
\begin{equation}
	U=\frac{U_0}{4d^2} \left(2z^2-x^2-y^2 \right) \eqtag{1}
\end{equation}
空间中存在磁场$\vec{B}=B_0\hat{z}$。本题不考虑可能的量子效应。

\subq{1}假设原点处有一个质量为$m$、电荷量大小为$q$的带正电的质点，固定电场不变，求解能使该质点局限在阱内做周期运动的$B_0$的范围。

\subq{2}实际的潘宁阱形成的电磁场不可能完美符合理论，答案只需保留一阶近似项，接下来分别考虑可能的误差：

\subsubq{2.1}考虑八极势带来的扰动，不考虑其它扰动，电势修正为
\begin{equation}
	U=\frac{U_0}{4d^2}\left(2z^2-x^2-y^2\right)+\frac{U_0}{2d^4}C\left[z^4-3z^2\left(x^2+y^2\right)+\frac{3}{8}\left(x^2+y^2\right)^2\right] \eqtag{2}
\end{equation}
其中$C\ll1$。已知粒子$z$方向运动振幅为$A$，求解粒子$z$方向运动基频相对变化量$\frac{\Delta \omega_z}{\omega_z}$，已知轴向运动主要受 
$\Delta U$中的$z^4$项影响（径向耦合项在时间平均后贡献很小，且频率不匹配）。

\subsubq{2.2}进一步考虑小量轴向阻尼力$-\gamma_z \dot{z}$，并施加一个正弦驱动力$\vec{F}=F_0\sin \left(\omega t+\theta\right) \hat{z}$进行驱动频率扫描，求解振幅关于频率的响应曲线在$\omega=\sqrt{\frac{U_0 q}{md^2}}$附近出现三值区间（即在此频率区间内，对应同一频率有三个不同的振幅解）的驱动力下界。

\subsubq{2.3}潘宁阱中粒子运动速率远小于光速，但当进行驱动频率扫描时，相对论效应会对振幅关于频率的响应曲线产生影响。此时粒子在阱中做圆周运动，驱动微波电场是右旋的圆极化波，频率为 $\omega$，振幅 $\mathcal{E}_0$，$\vec{\mathcal{E}}\left(t\right) = \mathcal{E}_0 \left(\hat{x}\cos\omega t - \hat{y}\sin\omega t \right)$，存在阻力$- \gamma_c \vec{v}$，求解频率$\omega$和运动动量大小$u$满足的关系，并说明保留适当的阶数时是否可能出现三值区间。简单起见，忽略空间电势的作用和粒子 $z$ 方向的运动。



\end{problemstatement}

% === 解答 ===
\begin{solution}

\solsubq{1}{10}
带电粒子在电磁场中的受力：
\begin{equation}
	\vec{F} = -q\nabla U -q\dot{x}B_0\hat{y}+q\dot{y}B_0\hat{x}\eqtagscore{1}{1} 
\end{equation}
先考虑$z$方向受力
\begin{equation}
	F_z=-\frac{U_0 q}{d^2} \eqtagscore{2}{2} z
\end{equation}
$z$ 方向被约束于阱内。考虑 $x$，$y$ 方向，设 $u=x+\mathrm{i} y$ ：
\begin{equation}
	m\ddot{u} = \frac{U_0 q}{2d^2}u-\mathrm{i} qB_0 \dot{u}\eqtagscore{3}{2} 
\end{equation}
令$u=A\mathrm{e}^{-\mathrm{i} \omega t}$：
\begin{equation}
	\omega^2-\frac{ q B_0}{m} \omega+ \frac{U_0 q}{2md^2} =0 \eqtagscore{4}{2}
\end{equation}
要求判别式大于$0$：
\begin{equation}
	B_0>\sqrt{\frac{2U_0 m}{qd^2}}\eqtagscore{5}{3}
\end{equation}

\solsubq{2}{30}
\solsubsubq{2.1}{6}
依题意写出略去径向影响的轴向运动方程：
\begin{equation}
	m\ddot{z}=-\frac{U_0 q}{d^2}z-\frac{2U_0 qC}{d^4}z^{3}\eqtagscore{6}{2}\label{eq:2.1}
\end{equation}
设 $z=A\sin \left(\omega_z t+\phi\right)$，代入式\ref{eq:2.1}并作小量近似：
\begin{equation}
	\omega_0=\sqrt{\frac{U_0q}{md^2}},\ \omega_z=\omega_0(1+\lambda C)\eqtag{7}
\end{equation}
利用三倍角公式化简，只保留一倍频：
\begin{equation}
	-\frac{U_0q}{d^2}(1+2\lambda C)A\sin(\omega_zt+\phi)=-\frac{U_0q}{d^2}A\sin(\omega_zt+\phi)-\frac{2U_0qC}{d^4}A^3\cdot\frac34\sin(\omega_zt+\phi)\eqtag{8}
\end{equation}
化简解得：
\begin{equation}
	\lambda=\dfrac{3A^2}{4d^2},\ \omega_z=\sqrt{\frac{U_0 q}{md^2}}\left(1+\frac{3CA^2}{4d^2}\right)\eqtagscore{9}{3}
\end{equation}
最终结果：
\begin{equation}
	\frac{\Delta \omega_z}{\omega_z} = \frac{3CA^2}{4d^2} \eqtagscore{10}{1}
\end{equation}

\solsubsubq{2.2}{12}
轴向运动方程：
\begin{equation}
	m\ddot{z}+\gamma_z \dot{z}+\frac{U_0 q}{d^2}z+\frac{2U_0 qC}{d^4}z^{3}=F_0\sin \left(\omega t+\theta\right) \eqtagscore{11}{2}\label{eq:2.2}
\end{equation}
将$z=A\sin \left(\omega t+\phi\right)$代入式\ref{eq:2.2}，利用三倍角公式化简，只保留一倍频：
\begin{equation}
	\begin{aligned}
		F_0\sin \left(\omega t+\theta\right)&=-mA\omega^2\sin \left(\omega t+\phi\right)+\gamma_z A\omega\cos \left(\omega t+\phi\right) \\
		&\quad +\frac{U_0 q}{d^2}A\sin \left(\omega t+\phi\right)+\frac{2U_0qC}{d^4}A^3\cdot\frac34\sin(\omega t+\phi)
	\end{aligned} \eqtagscore{12}{1}
\end{equation}
取模：
\begin{equation}
	\left(\frac{U_0 q}{d^2}-m\omega^2+\frac{3U_0 qCA^2}{2d^4}\right)^2 A^2+\gamma_z ^2 \omega^2 A^2=F_0^2\eqtagscore{13}{1}\label{eq:2.3}
\end{equation}
由$\frac{CA^2}{d^2} \ll 1$，将$\omega_z=\sqrt{\frac{U_0 q}{md^2}}\left(1+\frac{3CA^2}{4d^2}\right)$代入式\ref{eq:2.3}化简：
\begin{equation}
	\left(m\omega_z^2-m\omega^2\right)^2 A^2+\gamma_z ^2 \omega^2 A^2=F_0^2\eqtag{14}
\end{equation}
归一化处理，设 $w=\frac{\gamma_z^2A^2}{F_0^2}\omega_0^2,\ N=\frac{m(\omega_z-\omega_0)}{w\gamma_z},\ \delta=m\frac{\omega-\omega_0}{\gamma_z}$，整理得：
\begin{equation}
	\left(\frac{\omega+\omega_z}{\omega_0} \right)^2 \left(Nw-\delta\right)^2 w+ \left(\frac{\omega}{\omega_0} \right)^2 w=1 \eqtagscore{15}{1}
\end{equation}
此时 $\frac{\omega_0}{\omega}-1,\ \frac{\omega-\omega_z}{\omega_0}$应为小量，在小量近似下：
\begin{equation}
	w=\frac{1}{1+4\left(Nw-\delta\right)^2}\eqtagscore{16}{2} \label{eq:2.4}
\end{equation}
$A(\omega)$ 出现三值区间时，等效于 $w(\delta)$ 出现三值区间，因此 $\delta(w)$ 应当有至少一个稳定点，或以下方程应当有解：
\begin{equation}
	\frac{\mathrm{d}\delta}{\mathrm{d}w}=0\eqtagscore{17}{1}
\end{equation}
由式\ref{eq:2.4}解得$\delta$，计算得到：
\begin{equation}
	\delta=Nw \pm \sqrt{\frac{1}{4w}-\frac{1}{4}}\eqtag{18}\label{eq:2.5}
\end{equation}
要使 $\frac{\mathrm{d}\delta}{\mathrm{d}w}=0$ 可能有解，式\ref{eq:2.5}取$+$。
对式\ref{eq:2.5}两边求导：
\begin{equation}
	\frac{\mathrm{d}\delta}{\mathrm{d}w}=N-\frac{1}{4w^2\sqrt{\frac{1}{w}-1}}\eqtag{19}\label{eq:2.6}
\end{equation}
将 $\frac{\mathrm{d}\delta}{\mathrm{d}w}=0$ 代入式\ref{eq:2.6}，得到：
\begin{equation}
	N\geqslant\min\limits_{0<w<1}\frac{1}{4w^2\sqrt{\frac{1}{w}-1}}\eqtagscore{20}{1}
\end{equation}
右边最小值求得：
\begin{equation}
	N_{\min}=\frac{4\sqrt{3}}{9}\eqtagscore{21}{1}
\end{equation}

频率关于振幅的响应曲线出现三值区间的临界驱动力大小：
\begin{equation}
	F_0=\sqrt{\frac{16\sqrt{3}d\gamma_z^3}{27mC}\sqrt{\frac{U_0 q}{m}}}\eqtagscore{22}{2}
\end{equation}

\solsubsubq{2.3}{12}
相对论下运动方程：
\begin{equation}
	\frac{\mathrm{d}}{\mathrm{d}t}\left( \frac{\vec{v}}{\sqrt{1-\frac{|\vec{v}|^2}{c^2}}} \right) = -\frac{B_0 q}{m} \hat{z}\times\vec{v} - \frac{\gamma_c}{m} \vec{v} + \frac{q}{m}\vec{{\mathcal{E}}}\left(t\right)\eqtagscore{23}{2}
\end{equation}
用复数 $v=v_x+\mathrm{i}v_y$ 表示：
\begin{equation}
	\frac{\mathrm{d}}{\mathrm{d}t}\left( \frac{v}{\sqrt{1-\frac{|v|^2}{c^2}}} \right) =-\mathrm{i} \frac{B_0 q}{m} v - \frac{\gamma_c}{m} v + \frac{q}{m}\mathcal{E}_0 \mathrm{e}^{-\mathrm{i}\omega t}\eqtagscore{24}{2}
\end{equation}

引入旋转坐标系下的复动量 $u_c \mathrm{e}^{-\mathrm{i}\omega t} = \frac{mv}{\sqrt{1-\frac{|v|^2}{c^2}}}$，化简：
\begin{equation}
	\dot{u}_c - \mathrm{i}\left(\omega - \frac{B_0 q}{m\sqrt{1+\frac{|u_c|^2}{m^2c^2}}}\right) u_c + \frac{\gamma_c}{m\sqrt{1+\frac{|u_c|^2}{m^2c^2}}} u_c = q\mathcal{E}_0\eqtagscore{25}{3}
\end{equation}

稳定时 $\dot{u}_c=0$，取模：
\begin{equation}
	\left(\omega - \frac{B_0 q}{m\sqrt{1+\frac{u^2}{m^2c^2}}}\right) ^2 u^2 + \frac{\gamma_c ^2}{m^2\left(1+\frac{u^2}{m^2c^2}\right)} u^2 = q^2 \mathcal{E}_0 ^2\eqtagscore{26}{2}
\end{equation}
展开到一阶即有
\begin{equation}
	\left[\omega\left(1+\frac{u^2}{2m^2c^2}\right) - \frac{B_0 q}{m}\right] ^2 u^2 + \left(\frac{\gamma_c ^2}{m^2}-\frac{q^2\mathcal{E}_0^2}{m^2c^2}\right) u^2 = q^2 \mathcal{E}_0 ^2\eqtagscore{27}{1}
\end{equation}
与式\ref{eq:2.3}形式相似，可能出现三值区间。
\addtext{说明可能存在三值区间}{2}

注：若不忽略静电力，将得到
\begin{equation}
	\left[\omega\left(1+\frac{u^2}{2m^2c^2}\right) - \left(\frac{qU_0}{2d^2m\omega}+\frac{B_0 q}{m}\right)\right] ^2 u^2 + \left(\frac{\gamma_c ^2}{m^2}-\frac{q^2\mathcal{E}_0^2}{m^2c^2}\right) u^2 = q^2 \mathcal{E}_0 ^2\eqtag{28}
\end{equation}
也算作正确。



\scoring  % 评分标准自动使用题目环境中声明的总分
\end{solution}
	
\end{problem}

\end{document}
